% !TEX root = ../main.tex
\documentclass[../main.tex]{subfiles}
\graphicspath{{\subfix{../images/}}}

\begin{document}

Актуальность работы определяется необходимостью создания лёгких, переносимых и при этом достаточно выразительных игровых движков для обучения основам разработки компьютерных игр. Жанр <<bullet hell>> предъявляет повышенные требования к производительности и управлению большим количеством динамических объектов (снарядов), что делает его показательной задачей для отработки эффективных алгоритмов обработки данных, коллизий и событийного программирования. Библиотека Raylib, представляющая собой минималистичный, но полнофункциональный API для 2D- и 3D-графики, оконного управления и ввода, позволяет создавать игры без использования тяжёлых коммерческих движков (Unity, Unreal Engine) и при этом сохранять контроль над низкоуровневыми аспектами: циклами обновления, памятью, временем. Таким образом, разработка игры в жанре bullet hell на C++ с использованием Raylib является актуальной учебно-исследовательской задачей.

Новизна работы заключается в разработке событийно-ориентированной системы описания уровней, позволяющей гибко задавать временные последовательности атак (в том числе создавать вложенные события, когда одна атака порождает другие), а также в применении стандарта C++20 (в частности, \texttt{std::format} для форматирования строк интерфейса) и интеграции с кроссплатформенной системой сборки Premake/CMake. Такой подход упрощает добавление новых уровней и типов атак без изменения основного игрового цикла.

Цель работы: разработать на языке C++ с использованием библиотеки Raylib игру в жанре bullet hell, поддерживающую несколько уровней с различными паттернами атак и управление игроком с клавиатуры.

Объект исследования: методы организации игрового цикла, событийно-ориентированного планирования атак, обработки коллизий и отрисовки в реальном времени.

Предмет исследования: применение возможностей C++20 (включая стандартную библиотеку, контейнеры, алгоритмы, \texttt{std::format}) и библиотеки Raylib для реализации аркадной игры.

Задачи работы:
\begin{enumerate}
    \item Провести обзор существующих решений и выбрать технологический стек.
    \item Спроектировать архитектуру приложения: объектную модель игрока и снарядов, систему событий.
    \item Реализовать базовые механики: движение игрока (клавиши WASD/стрелки), нормализацию скорости, проверку столкновений (круг–прямоугольник), учёт неуязвимости после попадания.
    \item Разработать событийную систему для описания атак уровней, позволяющую создавать сложные последовательности (линии, лучи, взрывы, веера) и динамически добавлять новые события.
    \item Создать два демонстрационных уровня (обучающий и сложный) с использованием этой системы.
    \item Провести анализ полученных результатов.
\end{enumerate}

Практическая значимость работы заключается в том, что разработанный код может служить основой для дальнейшего развития (добавления новых типов атак, звукового сопровождения, редактора уровней) или использоваться в учебных целях при изучении C++ и игрового программирования.

\end{document}
